\chapter{Testing}
\label{chap:testing}

\defaultInstructions

\begin{length}
There are no strict length requirements for this chapter.  However, it is important that the writing is clear and concise.  Avoid repeating yourself.  Make sure to clearly signpost evidence for claims.  It is expected that a typical project needs 1--5 pages, but this can vary considerably from project to project.
\end{length}

\begin{expectations}
This chapter discusses the teams approach to software testing, as well as ensuring test quality.  Cover the automated test suite included with the source code, as well as any less visible forms testing the team employed to ensure the marking criteria were met.

If your team used manual testing as well as automated testing, it is important to include the following:
\begin{itemize}
\item What did your rely on manual testing for?  In other words, what was the purpose of manual testing.
\item To what extent did you rely on manual testing for aspects of software testing that would ideally be automated.
\item What are the respective limits of automated and manual testing.
\item An appendix listing each manual test.  For each manual test, describe precisely the actions the tester must undertake and what they are expected to observe.  Also add when the test was last performed and whether it was successful.
\end{itemize}
\end{expectations}